%% =====================================================================
%%  T-05-02  Follow the Interest
%%  -------------------------------------------------------------------
%%  One idea per file. Core is mandatory. Slots in canonical order.
%%  Max 700 words. No layout commands. No filler. New source material
%%  for this entry goes in sources/<BOOK>/T-05-02.tex -- never by
%%  rewriting this file (docs/RULES.md R0.1).
%% =====================================================================
%% @kind: topic
%% @id: T-05-02
%% @title: Follow the Interest
%% @subtitle: Watch where the benefit goes
%% @chapter: c05
%% @order: 20
%% @status: stable
%% @sources: B1, B3
%% @updated: 2026-09-19

\topic{T-05-02}{Follow the Interest}{Watch where the benefit goes}

\begin{core}
Ignore what people say they want and trace what they gain. In any arrangement the flow of benefit is visible long before the motive is stated, and it predicts behaviour with a reliability that no explanation matches.
\end{core}
\begin{mechanism}
Interests are structural: they follow from a person's position, their costs and their income, and they change slowly. Explanations are situational: they are produced for the audience, they change quickly, and they are chosen for effect. Tracing benefit works because it reads the structure rather than the presentation. It also works reflexively --- the fastest way to understand why someone is pressing you is to ask what your yes is worth to them, and the fastest way to be persuaded honestly is to ask what your no costs you.
\end{mechanism}
\begin{tells}
  \item The person proposing an arrangement describes your benefit at length and their own not at all.
  \item Terms favourable to you are argued for rather than simply stated.
  \item Someone's position changes when their incentives change, with no change in their reasoning.
  \item A third party with no visible stake is unusually invested in the outcome.
  \item The arrangement is complicated in a way that obscures who pays.
\end{tells}
\begin{moves}
  \item For every proposal, write down who gains, who pays, and when each happens.
  \item Check whether the gainer is the person explaining the arrangement.
  \item Trace the money or the credit through every intermediary before forming a view.
  \item Ask directly what the other party gets. The tone of the answer matters less than whether the question was expected.
  \item Assume a proposer has a gain --- nobody spends effort on an arrangement that does not benefit them.
\end{moves}
\begin{counters}
  \item Trace the benefit before evaluating the argument. An argument you cannot fault from a party who gains either way is not evidence of good faith.
  \item Require the terms in a form that shows who pays what. Complexity that survives a request to simplify is usually doing work for someone.
  \item Ask what the proposer would do if you said no. Their answer reveals the value of your yes.
  \item Be suspicious of arrangements where your benefit is asserted rather than demonstrated by the terms.
\end{counters}
\begin{cost}
Interest-tracing produces a cynical baseline, and the baseline is often wrong: people act from loyalty, confusion, principle and mistake at least as often as from gain. Used as a proof rather than as a question it generates confident misreadings and costs relationships. It is also weakest exactly where it feels strongest --- inside long relationships, where benefit flows are diffuse, mutual and poorly measured, and where the assumption of hidden gain does the most damage.
\end{cost}

\sources{B1, B3}
\seesources
\seealso{T-03-02, T-05-01, T-25-01}
